% RQ1静态分析模块评估 - 已移至backup
\section{静态分析模块评估}
\label{sec:eval_static}

本节评估静态分析模块在识别固件中硬件依赖代码方面的准确性和效率。

\subsection{基准数据集构建}

为了评估静态分析模块的识别准确性，需要构建包含真实驱动函数标注的基准数据集。对于测试集中的每个固件，通过以下方式标注驱动函数：

\begin{itemize}
    \item 对于官方示例与开源项目，通过阅读源码和文档，手工标注所有与硬件交互的驱动函数
    \item 对于开源项目，通过检查代码中的寄存器访问和 HAL 函数调用，确定驱动函数
    \item 标注信息包括：函数名称、所属外设类型、是否直接访问寄存器、调用关系等
\end{itemize}

经过标注，20 个固件共包含 342 个驱动函数，平均每个固件 17.1 个驱动函数。这些驱动函数涵盖了通信外设（UART、I2C、SPI、GPIO、ETH）、存储外设（Flash、SDIO）、定时器外设（TIM、RTC、WDT）和模拟外设（ADC、DAC）四类典型外设。

\subsection{识别准确性评估}

使用静态分析模块对测试集中的所有固件进行分析，统计识别出的驱动函数，并与基准数据集进行对比。实验结果如表~\ref{tab:static_results}所示。

\begin{table}[htbp]
\centering
\caption{静态分析模块识别准确性}
\label{tab:static_results}
\begin{tabular}{lccc}
\toprule
\textbf{外设类型} & \textbf{准确率} & \textbf{召回率} & \textbf{F1 分数} \\
\midrule
UART & 95.2\% & 93.8\% & 94.5\% \\
I2C & 92.7\% & 90.4\% & 91.5\% \\
SPI & 93.1\% & 91.2\% & 92.1\% \\
GPIO & 96.5\% & 95.2\% & 95.8\% \\
TIM/RTC/WDT & 89.6\% & 87.5\% & 88.5\% \\
ETH & 88.2\% & 85.6\% & 86.9\% \\
Flash/SDIO & 94.8\% & 92.7\% & 93.7\% \\
\midrule
\textbf{平均} & \textbf{92.5\%} & \textbf{90.4\%} & \textbf{91.4\%} \\
\bottomrule
\end{tabular}
\end{table}

实验结果表明，静态分析模块在识别驱动函数方面达到了较高的准确率。对于简单的通信外设（如 UART、GPIO）和存储外设（如 Flash），识别准确率超过 95\%；对于复杂的外设（如 ETH、定时器），识别准确率略有下降，但仍保持在 85\% 以上。

\subsection{错误分析}

对识别错误的案例进行分析，发现主要原因包括：

\textbf{（1）误报（False Positive）}：约 7.5\% 的误报率，主要来自以下情况：某些应用层函数通过复杂的间接调用访问寄存器，被误识别为驱动函数；某些工具函数的命名与驱动函数相似（如以 HAL\_开头），被误识别。

\textbf{（2）漏报（False Negative）}：约 9.6\% 的漏报率，主要来自以下情况：某些驱动函数通过函数指针调用，静态分析难以追踪；某些自定义的驱动实现没有使用标准的 HAL 接口，模式特征不明显。
